\section{Ample invertible sheaves}
\label{section-ample}

\noindent
Recall from Modules, Lemma \ref{modules-lemma-s-open}
that given an invertible sheaf $\mathcal{L}$ on a locally ringed
space $X$, and given a global section $s$ of $\mathcal{L}$
the set $X_s = \{x \in X \mid s \not \in \mathfrak m_x\mathcal{L}_x\}$
is open. A general remark is that
$X_s \cap X_{s'} = X_{ss'}$, where $ss'$ denote
the section $s \otimes s' \in \Gamma(X, \mathcal{L} \otimes \mathcal{L}')$.

\begin{definition}
\label{definition-ample}
\begin{reference}
\cite[II Definition 4.5.3]{EGA}
\end{reference}
Let $X$ be a scheme.
Let $\mathcal{L}$ be an invertible $\mathcal{O}_X$-module.
We say $\mathcal{L}$ is {\it ample} if
\begin{enumerate}
\item $X$ is quasi-compact, and
\item for every $x \in X$ there exists an $n \geq 1$
and $s \in \Gamma(X, \mathcal{L}^{\otimes n})$ such
that $x \in X_s$ and $X_s$ is affine.
\end{enumerate}
\end{definition}

\begin{lemma}
\label{lemma-ample-power-ample}
\begin{reference}
\cite[II Proposition 4.5.6(i)]{EGA}
\end{reference}
Let $X$ be a scheme. Let $\mathcal{L}$ be an invertible $\mathcal{O}_X$-module.
Let $n \geq 1$. Then $\mathcal{L}$ is ample if and only if
$\mathcal{L}^{\otimes n}$ is ample.
\end{lemma}

\begin{proof}
This follows from the fact that $X_{s^n} = X_s$.
\end{proof}

\begin{lemma}
\label{lemma-ample-on-closed}
Let $X$ be a scheme.
Let $\mathcal{L}$ be an ample invertible $\mathcal{O}_X$-module.
For any closed subscheme $Z \subset X$ the restriction of
$\mathcal{L}$ to $Z$ is ample.
\end{lemma}

\begin{proof}
This is clear since a closed subset of a quasi-compact space is quasi-compact
and a closed subscheme of an affine scheme is affine (see
Schemes, Lemma \ref{schemes-lemma-closed-immersion-affine-case}).
\end{proof}

\begin{lemma}
\label{lemma-affine-cap-s-open}
Let $X$ be a scheme. Let $\mathcal{L}$ be an invertible $\mathcal{O}_X$-module.
Let $s \in \Gamma(X, \mathcal{L})$. For any affine $U \subset X$
the intersection $U \cap X_s$ is affine.
\end{lemma}

\begin{proof}
This translates into the following algebra problem.
Let $R$ be a ring. Let $N$ be an invertible $R$-module
(i.e., locally free of rank 1). Let $s \in N$ be an element.
Then $V = \{\mathfrak p \mid s \not \in \mathfrak p N\}$ is
an affine open subset of $\Spec(R)$.

\medskip\noindent
Let $A = \bigoplus_{n \geq 0} A_n$ be the symmetric algebra of $N$
(which is commutative) and view $s$ as an element of $A_1$.
Set $B = A/(s - 1)A$. This is an $R$-algebra whose construction
commutes with any base change $R \to R'$. Thus $B' = B \otimes_R R'$
is the zero ring if $s$ maps to zero in $N' = N \otimes_R R'$. It
follows that if $x \in \Spec(R) \setminus V$, then $B \otimes_R \kappa(x) = 0$.
We conclude that $\Spec(B) \to \Spec(R)$ factors through $V$ as the
fibres over $x \not \in V$ are empty. On the other hand, if
$\Spec(R') \subset V$ is an affine open, then $s$ maps to a basis
element of $N'$ and we see that $B' = R'[s]/(s - 1) \cong R'$.
It follows that $\Spec(B) \to V$ is an isomorphism and $V$ is
indeed affine.
\end{proof}

\begin{lemma}
\label{lemma-ample-tensor-globally-generated}
\begin{reference}
\cite[II Proposition 4.5.6(ii)]{EGA}
\end{reference}
Let $X$ be a scheme. Let $\mathcal{L}$ and $\mathcal{M}$
be invertible $\mathcal{O}_X$-modules. If
\begin{enumerate}
\item $\mathcal{L}$ is ample, and
\item the open sets $X_t$ where $t \in \Gamma(X, \mathcal{M}^{\otimes m})$
for $m > 0$ cover $X$,
\end{enumerate}
then $\mathcal{L} \otimes \mathcal{M}$ is ample.
\end{lemma}

\begin{proof}
We check the conditions of Definition \ref{definition-ample}.
As $\mathcal{L}$ is ample we see that $X$ is quasi-compact.
Let $x \in X$. Choose $n \geq 1$, $m \geq 1$,
$s \in \Gamma(X, \mathcal{L}^{\otimes n})$, and
$t \in \Gamma(X, \mathcal{M}^{\otimes m})$
such that $x \in X_s$, $x \in X_t$ and $X_s$ is affine.
Then $s^mt^n \in \Gamma(X, (\mathcal{L} \otimes \mathcal{M})^{\otimes nm})$,
$x \in X_{s^mt^n}$, and $X_{s^mt^n}$ is affine by
Lemma \ref{lemma-affine-cap-s-open}.
\end{proof}

\begin{lemma}
\label{lemma-affine-s-opens}
Let $X$ be a scheme. Let $\mathcal{L}$ be an invertible $\mathcal{O}_X$-module.
Assume the open sets $X_s$, where $s \in \Gamma(X, \mathcal{L}^{\otimes n})$
and $n \geq 1$, form a basis for the topology on $X$.
Then among those opens, the open sets $X_s$ which are affine
form a basis for the topology on $X$.
\end{lemma}

\begin{proof}
Let $x \in X$. Choose an affine open neighbourhood
$\Spec(R) = U \subset X$ of $x$.
By assumption, there exists
a $n \geq 1$ and a $s \in \Gamma(X, \mathcal{L}^{\otimes n})$
such that $X_s \subset U$. By Lemma \ref{lemma-affine-cap-s-open} above
the intersection $X_s = U \cap X_s$ is affine.
Since $U$ can be chosen arbitrarily small we win.
\end{proof}

\begin{lemma}
\label{lemma-affine-s-opens-cover-quasi-separated}
Let $X$ be a scheme and $\mathcal{L}$ be an invertible $\mathcal{O}_X$-module.
Assume for every point $x$ of $X$ there exists $n \geq 1$ and
$s \in \Gamma(X, \mathcal{L}^{\otimes n})$ such that
$x \in X_s$ and $X_s$ is affine. Then $X$ is separated.
\end{lemma}

\begin{proof}
We show first that $X$ is quasi-separated.  By assumption we
can find a covering of $X$ by affine opens of the form
$X_s$.  By Lemma \ref{lemma-affine-cap-s-open}, the
intersection of any two such sets is affine, so Schemes, Lemma
\ref{schemes-lemma-characterize-quasi-separated} implies
that $X$ is quasi-separated.

\medskip\noindent
To show that $X$ is separated, we can use the valuative
criterion, Schemes, Lemma
\ref{schemes-lemma-valuative-criterion-separatedness}.
Thus, let $A$ be a valuation ring with fraction field $K$
and consider two morphisms $f, g : \Spec(A) \to X$ such that
the two compositions $\Spec(K) \to \Spec(A) \to X$ agree.
As $A$ is local, there exists $p, q \ge 1$, $s \in \Gamma(X,
\mathcal{L}^{\otimes p})$, and $t \in \Gamma(X,
\mathcal{L}^{\otimes q})$ such that $X_s$ and $X_t$ are
affine, $f(\Spec A) \subseteq X_s$, and $g(\Spec A)
\subseteq X_t$.  We now replace $s$ by $s^q$, $t$ by $t^p$,
and $\mathcal{L}$ by $\mathcal{L}^{\otimes pq}$.  This is
harmless as $X_s = X_{s^q}$ and $X_t = X_{t^p}$, and now $s$
and $t$ are both sections of the same sheaf $\mathcal{L}$.

\medskip\noindent
The quasi-coherent module $f^*\mathcal{L}$ corresponds to an $A$-module $M$ and
$g^*\mathcal{L}$ corresponds to an $A$-module $N$ by our
classification of quasi-coherent modules over affine schemes
(Schemes, Lemma \ref{schemes-lemma-quasi-coherent-affine}).
The $A$-modules $M$ and $N$ are locally free of rank
$1$ (Lemma \ref{lemma-locally-free-module}) and as $A$ is
local they are free (Algebra, Lemma
\ref{algebra-lemma-K0-local}).  Therefore we may identify
$M$ and $N$ with $A$-submodules of $M \otimes_A K$ and $N
\otimes_A K$.  The equality $f|_{\Spec(K)} = g|_{\Spec(K)}$
determines an isomorphism $\phi \colon M \otimes_A K \to N
\otimes_A K$.

\medskip\noindent
Let $x \in M$ and $y \in N$ be the elements corresponding to
the pullback of $s$ along $f$ and $g$, respectively.  These
satisfy $\phi(x \otimes 1) = y \otimes 1$.  The image of $f$
is contained in $X_s$, so $x \not\in \mathfrak{m}_A M$, that
is, $x$ generates $M$.  Hence $\phi$ determines an
isomorphism of $M$ with the submodule of $N$ generated by
$y$.  Arguing symmetrically using $t$, $\phi^{-1}$
determines an isomorphism of $N$ with a submodule of $M$.
Consequently $\phi$ restricts to an isomorphism of $M$ and
$N$.  Since $x$ generates $M$, its image $y$ generates
$N$, implying $y \not\in \mathfrak{m}_A N$.  Therefore
$g(\Spec(A)) \subseteq X_s$.  Because $X_s$ is affine, it is
separated by Schemes, Lemma \ref{schemes-lemma-affine-separated},
and we conclude $f = g$.
\end{proof}

\begin{lemma}
\label{lemma-ample-separated}
Let $X$ be a scheme. If there exists an ample invertible sheaf on $X$
then $X$ is separated.
\end{lemma}

\begin{proof}
Follows immediately from
Lemma \ref{lemma-affine-s-opens-cover-quasi-separated} and
Definition \ref{definition-ample}.
\end{proof}

\begin{lemma}
\label{lemma-map-into-proj}
Let $X$ be a scheme.
Let $\mathcal{L}$ be an invertible $\mathcal{O}_X$-module.
Set $S = \Gamma_*(X, \mathcal{L})$ as a graded ring.
If every point of $X$ is contained in one of the
open subschemes $X_s$, for some $s \in S_{+}$ homogeneous, then
there is a canonical morphism of schemes
$$
f : X \longrightarrow Y = \text{Proj}(S),
$$
to the homogeneous spectrum of $S$ (see
Constructions, Section \ref{constructions-section-proj}).
This morphism has the following properties
\begin{enumerate}
\item $f^{-1}(D_{+}(s)) = X_s$ for any $s \in S_{+}$ homogeneous,
\item there are $\mathcal{O}_X$-module maps
$f^*\mathcal{O}_Y(n) \to \mathcal{L}^{\otimes n}$
compatible with multiplication maps, see
Constructions, Equation (\ref{constructions-equation-multiply}),
\item the composition
$S_n \to \Gamma(Y, \mathcal{O}_Y(n)) \to \Gamma(X, \mathcal{L}^{\otimes n})$
is the identity map, and
\item for every $x \in X$ there is an integer $d \geq 1$
and an open neighbourhood $U \subset X$ of $x$
such that $f^*\mathcal{O}_Y(dn)|_U \to \mathcal{L}^{\otimes dn}|_U$
is an isomorphism for all $n \in \mathbf{Z}$.
\end{enumerate}
\end{lemma}

\begin{proof}
Denote $\psi : S \to \Gamma_*(X, \mathcal{L})$ the identity map.
We are going to use the triple
$(U(\psi), r_{\mathcal{L}, \psi}, \theta)$ of
Constructions, Lemma \ref{constructions-lemma-invertible-map-into-proj}.
By assumption the open subscheme $U(\psi)$ of equals $X$. Hence
$r_{\mathcal{L}, \psi} : U(\psi) \to Y$ is defined on all of $X$.
We set $f = r_{\mathcal{L}, \psi}$.
The maps in part (2) are the components of $\theta$.
Part (3) follows from condition (2) in the lemma cited above.
Part (1) follows from (3) combined with condition (1) in the lemma
cited above. Part (4) follows from the last statement in
Constructions, Lemma \ref{constructions-lemma-invertible-map-into-proj}
since the map $\alpha$ mentioned there is an isomorphism.
\end{proof}

\begin{lemma}
\label{lemma-map-into-proj-quasi-compact}
Let $X$ be a scheme. Let $\mathcal{L}$ be an invertible $\mathcal{O}_X$-module.
Set $S = \Gamma_*(X, \mathcal{L})$.
Assume (a) every point of $X$ is contained in one of the
open subschemes $X_s$, for some $s \in S_{+}$ homogeneous,
and (b) $X$ is quasi-compact. Then the canonical morphism of schemes
$f : X \longrightarrow \text{Proj}(S)$ of Lemma \ref{lemma-map-into-proj}
above is quasi-compact with dense image.
\end{lemma}

\begin{proof}
To prove $f$ is quasi-compact it suffices to show that $f^{-1}(D_{+}(s))$
is quasi-compact for any $s \in S_{+}$ homogeneous. Write
$X = \bigcup_{i = 1, \ldots, n} X_i$ as a finite union of
affine opens. By Lemma \ref{lemma-affine-cap-s-open} each intersection
$X_s \cap X_i$ is affine. Hence $X_s = \bigcup_{i = 1, \ldots, n} X_s \cap X_i$
is quasi-compact. Assume that the image of $f$ is not dense to get
a contradiction. Then, since the opens $D_+(s)$ with $s \in S_+$ homogeneous
form a basis for the topology on $\text{Proj}(S)$, we can find such
an $s$ with $D_+(s) \not = \emptyset$ and $f(X) \cap D_+(s) = \emptyset$.
By Lemma \ref{lemma-map-into-proj}
this means $X_s = \emptyset$. By Lemma \ref{lemma-invert-s-sections}
this means that a power $s^n$ is the zero section of
$\mathcal{L}^{\otimes n\deg(s)}$.
This in turn means that $D_+(s) = \emptyset$ which is the
desired contradiction.
\end{proof}

\begin{lemma}
\label{lemma-ample-immersion-into-proj}
Let $X$ be a scheme. Let $\mathcal{L}$ be an invertible $\mathcal{O}_X$-module.
Set $S = \Gamma_*(X, \mathcal{L})$.
Assume $\mathcal{L}$ is ample. Then the canonical morphism of schemes
$f : X \longrightarrow \text{Proj}(S)$ of Lemma \ref{lemma-map-into-proj}
is an open immersion with dense image.
\end{lemma}

\begin{proof}
By Lemma \ref{lemma-affine-s-opens-cover-quasi-separated} we see
that $X$ is quasi-separated. Choose finitely many
$s_1, \ldots, s_n \in S_{+}$ homogeneous
such that $X_{s_i}$ are affine, and $X = \bigcup X_{s_i}$.
Say $s_i$ has degree $d_i$. The inverse image of
$D_{+}(s_i)$ under $f$ is $X_{s_i}$, see Lemma \ref{lemma-map-into-proj}.
By Lemma \ref{lemma-invert-s-sections} the ring map
$$
(S^{(d_i)})_{(s_i)} = \Gamma(D_{+}(s_i), \mathcal{O}_{\text{Proj}(S)})
\longrightarrow
\Gamma(X_{s_i}, \mathcal{O}_X)
$$
is an isomorphism. Hence $f$ induces an isomorphism
$X_{s_i} \to D_{+}(s_i)$. Thus $f$ is an isomorphism of $X$ onto the open
subscheme $\bigcup_{i = 1, \ldots, n} D_{+}(s_i)$ of $\text{Proj}(S)$.
The image is dense by Lemma \ref{lemma-map-into-proj-quasi-compact}.
\end{proof}

\begin{lemma}
\label{lemma-open-in-proj-ample}
Let $X$ be a scheme.
Let $S$ be a graded ring. Assume $X$ is quasi-compact,
and assume there exists an open immersion
$$
j : X \longrightarrow Y = \text{Proj}(S).
$$
Then $j^*\mathcal{O}_Y(d)$ is an invertible ample sheaf
for some $d > 0$.
\end{lemma}

\begin{proof}
This is Constructions, Lemma \ref{constructions-lemma-ample-on-proj}.
\end{proof}

\begin{proposition}
\label{proposition-characterize-ample}
Let $X$ be a quasi-compact scheme.
Let $\mathcal{L}$ be an invertible sheaf on $X$.
Set $S = \Gamma_*(X, \mathcal{L})$.
The following are equivalent:
\begin{enumerate}
\item
\label{item-ample}
$\mathcal{L}$ is ample,
\item
\label{item-immersion}
the open sets $X_s$, with $s \in S_{+}$ homogeneous,
cover $X$ and the associated morphism $X \to \text{Proj}(S)$
is an open immersion,
\item
\label{item-s-basis}
the open sets $X_s$, with $s \in S_{+}$ homogeneous,
form a basis for the topology of $X$,
\item
\label{item-s-affine-basis}
the open sets $X_s$, with $s \in S_{+}$ homogeneous,
which are affine form a basis for the topology of $X$,
\item
\label{item-qc-gg}
for every quasi-coherent sheaf $\mathcal{F}$ on $X$
the sum of the images of the canonical maps
$$
\Gamma(X, \mathcal{F} \otimes_{\mathcal{O}_X} \mathcal{L}^{\otimes n})
\otimes_{\mathbf{Z}} \mathcal{L}^{\otimes -n}
\longrightarrow
\mathcal{F}
$$
with $n \geq 1$ equals $\mathcal{F}$,
\item
\label{item-qc-i-gg}
same property as (\ref{item-qc-gg}) with $\mathcal{F}$
ranging over all quasi-coherent sheaves of ideals,
\item
\label{item-c-gg}
$X$ is quasi-separated and
for every quasi-coherent sheaf $\mathcal{F}$ of finite type on $X$
there exists an integer $n_0$ such that
$\mathcal{F} \otimes_{\mathcal{O}_X} \mathcal{L}^{\otimes n}$
is globally generated for all $n \geq n_0$,
\item
\label{item-c-q}
$X$ is quasi-separated and
for every quasi-coherent sheaf $\mathcal{F}$ of finite type on $X$
there exist integers $n > 0$, $k \geq 0$ such that
$\mathcal{F}$ is a quotient of a direct sum of $k$ copies of
$\mathcal{L}^{\otimes - n}$, and
\item
\label{item-c-i-q}
same as in (\ref{item-c-q}) with $\mathcal{F}$ ranging over all
sheaves of ideals of finite type on $X$.
\end{enumerate}
\end{proposition}

\begin{proof}
Lemma \ref{lemma-ample-immersion-into-proj} is
(\ref{item-ample}) $\Rightarrow$ (\ref{item-immersion}).
Lemmas \ref{lemma-ample-power-ample} and \ref{lemma-open-in-proj-ample}
provide the implication
(\ref{item-ample}) $\Leftarrow$ (\ref{item-immersion}).
The implications (\ref{item-immersion}) $\Rightarrow$
(\ref{item-s-affine-basis}) $\Rightarrow$ (\ref{item-s-basis})
are clear from Constructions, Section \ref{constructions-section-proj}.
Lemma \ref{lemma-affine-s-opens} is
(\ref{item-s-basis}) $\Rightarrow$ (\ref{item-ample}).
Thus we see that the first 4 conditions are all equivalent.

\medskip\noindent
Assume the equivalent conditions (1) -- (4).
Note that in particular $X$ is separated (as an open
subscheme of the separated scheme $\text{Proj}(S)$).
Let $\mathcal{F}$ be a quasi-coherent sheaf on $X$.
Choose $s \in S_{+}$ homogeneous such that $X_s$ is affine.
We claim that any section $m \in \Gamma(X_s, \mathcal{F})$
is in the image of one of the maps displayed in
(\ref{item-qc-gg}) above. This will imply (\ref{item-qc-gg})
since these affines $X_s$ cover $X$.
Namely, by Lemma \ref{lemma-invert-s-sections} we may write
$m$ as the image of $m' \otimes s^{-n}$ for some
$n \geq 1$, some
$m' \in \Gamma(X, \mathcal{F} \otimes \mathcal{L}^{\otimes n})$.
This proves the claim.

\medskip\noindent
Clearly (\ref{item-qc-gg}) $\Rightarrow$ (\ref{item-qc-i-gg}).
Let us assume (\ref{item-qc-i-gg}) and prove $\mathcal{L}$ is
ample. Pick $x \in X$. Let $U \subset X$ be an affine open
which contains $x$. Set $Z = X \setminus U$. We may think of
$Z$ as a reduced closed subscheme, see
Schemes, Section \ref{schemes-section-reduced}.
Let $\mathcal{I} \subset \mathcal{O}_X$ be the quasi-coherent
sheaf of ideals corresponding to the closed subscheme $Z$.
By assumption (\ref{item-qc-i-gg}), there exists an $n \geq 1$ and a section
$s \in \Gamma(X, \mathcal{I} \otimes \mathcal{L}^{\otimes n})$
such that $s$ does not vanish at $x$ (more precisely such that
$s \not \in \mathfrak m_x \mathcal{I}_x \otimes \mathcal{L}_x^{\otimes n}$).
We may think of $s$ as a section of $\mathcal{L}^{\otimes n}$.
Since it clearly vanishes along $Z$ we see that
$X_s \subset U$. Hence $X_s$ is affine, see
Lemma \ref{lemma-affine-cap-s-open}.
This proves that $\mathcal{L}$ is ample.
At this point we have proved that (1) -- (6) are equivalent.

\medskip\noindent
Assume the equivalent conditions (1) -- (6). In the following
we will use the fact that the tensor product of two sheaves of
modules which are globally generated is globally generated without
further mention (see
Modules, Lemma \ref{modules-lemma-tensor-product-globally-generated}).
By (1) we can find elements $s_i \in S_{d_i}$ with $d_i \geq 1$
such that $X = \bigcup_{i = 1, \ldots, n} X_{s_i}$.
Set $d = d_1\ldots d_n$. It follows that $\mathcal{L}^{\otimes d}$
is globally generated by
$$
s_1^{d/d_1}, \ldots, s_n^{d/d_n}.
$$
This means that if $\mathcal{L}^{\otimes j}$ is globally generated
then so is $\mathcal{L}^{\otimes j + dn}$ for all $n \geq 0$.
Fix a $j \in \{0, \ldots, d - 1\}$. For any point $x \in X$ there
exists an $n \geq 1$ and a global section $s$ of $\mathcal{L}^{j + dn}$
which does not vanish at $x$, as follows from (\ref{item-qc-gg}) applied
to $\mathcal{F} = \mathcal{L}^{\otimes j}$ and ample invertible
sheaf $\mathcal{L}^{\otimes d}$. Since $X$ is quasi-compact there
we may find a finite list of integers $n_i$ and global sections
$s_i$ of $\mathcal{L}^{\otimes j + dn_i}$ which do not vanish at any point
of $X$. Since $\mathcal{L}^{\otimes d}$ is globally generated this means that
$\mathcal{L}^{\otimes j + dn}$ is globally generated where $n = \max\{n_i\}$.
Since we proved this for every congruence class mod $d$ we
conclude that there exists an $n_0 = n_0(\mathcal{L})$ such that
$\mathcal{L}^{\otimes n}$ is globally generated for all $n \geq n_0$.
At this point we see that if $\mathcal{F}$ is globally generated then
so is $\mathcal{F} \otimes \mathcal{L}^{\otimes n}$ for all
$n \geq n_0$.

\medskip\noindent
We continue to assume the equivalent conditions (1) -- (6).
Let $\mathcal{F}$ be a quasi-coherent
sheaf of $\mathcal{O}_X$-modules of finite type.
Denote $\mathcal{F}_n \subset \mathcal{F}$ the image of the canonical
map of (\ref{item-qc-gg}). By construction
$\mathcal{F}_n \otimes \mathcal{L}^{\otimes n}$ is
globally generated. By (\ref{item-qc-gg}) we see
$\mathcal{F}$ is the sum of the subsheaves $\mathcal{F}_n$,
$n \geq 1$. By
Modules, Lemma \ref{modules-lemma-finite-type-quasi-compact-colimit}
we see that $\mathcal{F} = \sum_{n = 1, \ldots, N} \mathcal{F}_n$
for some $N \geq 1$. It follows that
$\mathcal{F} \otimes \mathcal{L}^{\otimes n}$ is globally
generated whenever $n \geq N + n_0(\mathcal{L})$ with $n_0(\mathcal{L})$
as above. We conclude that (1) -- (6) implies (\ref{item-c-gg}).

\medskip\noindent
Assume (\ref{item-c-gg}). Let $\mathcal{F}$ be a quasi-coherent
sheaf of $\mathcal{O}_X$-modules of finite type.
By (\ref{item-c-gg}) there exists an integer $n \geq 1$ such that
the canonical map
$$
\Gamma(X, \mathcal{F} \otimes_{\mathcal{O}_X} \mathcal{L}^{\otimes n})
\otimes_{\mathbf{Z}} \mathcal{L}^{\otimes -n}
\longrightarrow
\mathcal{F}
$$
is surjective. Let $I$ be the set of finite subsets of
$\Gamma(X, \mathcal{F} \otimes_{\mathcal{O}_X} \mathcal{L}^{\otimes n})$
partially ordered by inclusion. Then $I$ is a directed partially ordered set.
For $i = \{s_1, \ldots, s_{r(i)}\}$ let $\mathcal{F}_i \subset \mathcal{F}$
be the image of the map
$$
\bigoplus\nolimits_{j = 1, \ldots, r(i)} \mathcal{L}^{\otimes -n}
\longrightarrow
\mathcal{F}
$$
which is multiplication by $s_j$ on the $j$th factor. The surjectivity above
implies that $\mathcal{F} = \colim_{i \in I} \mathcal{F}_i$.
Hence Modules, Lemma \ref{modules-lemma-finite-type-quasi-compact-colimit}
applies and we conclude that
$\mathcal{F} = \mathcal{F}_i$ for some $i$.
Hence we have proved (\ref{item-c-q}). In other words,
(\ref{item-c-gg}) $\Rightarrow$ (\ref{item-c-q}).

\medskip\noindent
The implication (\ref{item-c-q}) $\Rightarrow$ (\ref{item-c-i-q}) is trivial.

\medskip\noindent
Finally, assume (\ref{item-c-i-q}).
Let $\mathcal{I} \subset \mathcal{O}_X$ be a quasi-coherent sheaf
of ideals. By Lemma \ref{lemma-quasi-coherent-colimit-finite-type}
(this is where we use the condition that $X$ be quasi-separated)
we see that $\mathcal{I} = \colim_\alpha I_\alpha$ with
each $I_\alpha$ quasi-coherent of finite type. Since by assumption each of
the $I_\alpha$ is a quotient of negative tensor powers of
$\mathcal{L}$ we conclude the same for $\mathcal{I}$ (but of course
without the finiteness or boundedness of the powers). Hence
we conclude that (\ref{item-c-i-q}) implies (\ref{item-qc-i-gg}).
This ends the proof of the proposition.
\end{proof}

\begin{lemma}
\label{lemma-ample-on-locally-closed}
Let $X$ be a scheme. Let $\mathcal{L}$ be an ample invertible
$\mathcal{O}_X$-module. Let $i : X' \to X$ be a morphism of schemes.
Assume at least one of the following conditions holds
\begin{enumerate}
\item $i$ is a quasi-compact immersion,
\item $X'$ is quasi-compact and $i$ is an immersion,
\item $i$ is quasi-compact and induces a homeomorphism
between $X'$ and $i(X')$,
\item $X'$ is quasi-compact and $i$ induces a homeomorphism
between $X'$ and $i(X')$.
\end{enumerate}
Then $i^*\mathcal{L}$ is ample on $X'$.
\end{lemma}

\begin{proof}
Observe that in cases (1) and (3) the scheme $X'$ is quasi-compact
as $X$ is quasi-compact by Definition \ref{definition-ample}.
Thus it suffices to prove (2) and (4). Since (2) is a special case
of (4) it suffices to prove (4).

\medskip\noindent
Assume condition (4) holds. For $s \in \Gamma(X, \mathcal{L}^{\otimes d})$
denote $s' = i^*s$ the pullback of $s$ to $X'$. Note that
$s'$ is a section of $(i^*\mathcal{L})^{\otimes d}$. By
Proposition \ref{proposition-characterize-ample}
the opens $X_s$, for $s \in \Gamma(X, \mathcal{L}^{\otimes d})$,
form a basis for the topology on $X$. Since $X'_{s'} = i^{-1}(X_s)$ by
Modules, Remark \ref{modules-remark-pullback-s-open}
and since $X' \to i(X')$ is a homeomorphism, we conclude
the opens $X'_{s'}$ form a basis for the topology of $X'$. Hence
$i^*\mathcal{L}$ is ample by
Proposition \ref{proposition-characterize-ample}.
\end{proof}

\begin{lemma}
\label{lemma-ample-on-product}
Let $S$ be a quasi-separated scheme. Let $X$, $Y$ be schemes over $S$.
Let $\mathcal{L}$ be an ample invertible $\mathcal{O}_X$-module
and let $\mathcal{N}$ be an ample invertible $\mathcal{O}_Y$-module.
Then $\mathcal{M} = \text{pr}_1^*\mathcal{L}
\otimes_{\mathcal{O}_{X \times_S Y}} \text{pr}_2^*\mathcal{N}$
is an ample invertible sheaf on $X \times_S Y$.
\end{lemma}

\begin{proof}
The morphism $i : X \times_S Y \to X \times Y$ is a quasi-compact
immersion, see Schemes, Lemma \ref{schemes-lemma-fibre-product-after-map}.
On the other hand, $\mathcal{M}$ is the pullback by
$i$ of the corresponding invertible module on $X \times Y$.
By Lemma \ref{lemma-ample-on-locally-closed} it suffices to prove the
lemma for $X \times Y$. We check (1) and (2) of
Definition \ref{definition-ample} for $\mathcal{M}$ on $X \times Y$.

\medskip\noindent
Since $X$ and $Y$ are quasi-compact, so is $X \times Y$.
Let $z \in X \times Y$ be a point. Let $x \in X$ and $y \in Y$
be the projections. Choose $n > 0$ and
$s \in \Gamma(X, \mathcal{L}^{\otimes n})$
such that $X_s$ is an affine open neighbourhood of $x$.
Choose $m > 0$ and
$t \in \Gamma(Y, \mathcal{N}^{\otimes m})$
such that $Y_t$ is an affine open neighbourhood of $y$.
Then $r = \text{pr}_1^*s \otimes \text{pr}_2^*t$ is a section
of $\mathcal{M}$ with $(X \times Y)_r = X_s \times Y_t$.
This is an affine open neighbourhood of $z$ and the proof is complete.
\end{proof}












